\section*{TÀI LIỆU THAM KHẢO}
\addcontentsline{toc}{section}{TÀI LIỆU THAM KHẢO}

\begin{enumerate}
    \item Đào Duy Anh (2014), \textit{Việt Nam văn hóa sử cương}, Quan Hải tùng thư xuất bản năm 1938; Nhà xuất bản Thế giới -- Công ty Văn hóa và Truyền thông Nhã Nam tái bản, Hà Nội.
    \item Ban Chấp hành Trung ương Đảng Cộng sản Việt Nam (2023), \textit{Nghị quyết số 45-NQ/TW ngày 24/11/2023 của Hội nghị lần thứ tám Ban Chấp hành Trung ương Đảng Khóa XIII về tiếp tục xây dựng và phát huy vai trò của đội ngũ trí thức đáp ứng yêu cầu phát triển đất nước nhanh và bền vững trong giai đoạn mới}, Hà Nội.
    \item Bộ Giáo dục và Đào tạo (2021), \textit{Giáo trình Triết học Mác -- Lênin} (Dành cho bậc đào tạo đại học và sau đại học hệ không chuyên lý luận chính trị), Nhà xuất bản Chính trị quốc gia Sự thật, Hà Nội.
    \item Lê Quý Đôn (1977), \textit{Vân đài loại ngữ} (Bản dịch của Viện Sử học), Nhà xuất bản Văn hóa Thông tin, Hà Nội.
    \item Lê Quý Đôn (1977), \textit{Phủ biên tạp lục} (Bản dịch của Viện Sử học), Nhà xuất bản Khoa học Xã hội, Hà Nội.
    \item Trần Văn Giàu (1973), \textit{Sự phát triển của tư tưởng ở Việt Nam từ thế kỷ XIX đến Cách mạng Tháng Tám} (Tập 1: Hệ ý thức phong kiến và sự thất bại của nó trước các nhiệm vụ lịch sử), Nhà xuất bản Khoa học Xã hội, Hà Nội.
    \item Phan Huy Lê, Chu Thiên, Vương Hoàng Tuyên, Đinh Xuân Lâm (1965), \textit{Lịch sử chế độ phong kiến Việt Nam} (Tập III -- Thời kỳ khủng hoảng và phong trào nông dân khởi nghĩa thế kỷ XVIII), Nhà xuất bản Giáo dục, Hà Nội.
    \item Ngô Thì Nhậm (2001), \textit{Trúc Lâm tông chỉ nguyên thanh} (Bản dịch của Viện Nghiên cứu Hán Nôm), Nhà xuất bản Hà Nội, Hà Nội.
    \item Nguyễn Khắc Thuần (1998), \textit{Đại cương lịch sử văn hóa Việt Nam} (Tập 4), Nhà xuất bản Giáo dục, Hà Nội.
    \item Nguyễn Tài Thư (Chủ biên, 1993), \textit{Lịch sử tư tưởng Việt Nam} (Tập 1), Viện Triết học -- Viện Khoa học Xã hội Việt Nam, Nhà xuất bản Khoa học Xã hội, Hà Nội.
    \item Đặng Việt Bích (2015), ``Quang Trung trong tâm trí kẻ sĩ Bắc Hà'', \textit{Tạp chí Văn học Nghệ thuật}.
    \item Quách Hiền (2018), ``Tư học Việt Nam thế kỷ XVIII'' (Khảo sát quá trình chuyển dịch của các danh nho về làng xã mở trường tư thục), chuyên trang Nghiên Cứu Lịch Sử.
    \item Nguyễn Thị Hồng Phương (2026), ``Nhân sinh quan trong tư tưởng triết học của Lê Quý Đôn -- Ý nghĩa đối với xây dựng con người Việt Nam hiện nay'', \textit{Tạp chí Nghiên cứu Triết học}, Số 31 (2-2026), tr. 62--69.
    \item Nguyễn Kim Sơn (2004), ``Những nhân tố làm khởi phát khuynh hướng khảo chứng học thế kỷ XVIII -- nửa đầu thế kỷ XIX'', \textit{Kỷ yếu Thông báo Hán Nôm học}, Viện Nghiên cứu Hán Nôm -- Viện Hàn lâm Khoa học xã hội Việt Nam.
    \item Cao Văn Thức (2016), ``Tản mạn về học vị Sinh Đồ và chuyện Sinh Đồ ba quan'', chuyên trang Nghiên Cứu Lịch Sử.
    \item Sở Giáo dục và Đào tạo Thành phố Hồ Chí Minh và Thư viện tỉnh Đồng Nai (2020), Chuyên đề: ``Danh sư Võ Trường Toản -- Bậc Sùng Đức xử sĩ và hào khí Đồng Nai xứ Gia Định thế kỷ XVIII''.
\end{enumerate}

